\scanpage{13}

Sous ces conditions, \unsure{la|\deleted{la}} factorisation (i) est unique à isom près, et $S'$ est de \deleted{t.f.}\inserted{prés. finie} sur $S$.

(i) $\Rightarrow$ (i bis) $\Rightarrow$ (ii) trivial. De $=$ l'unicité de la factorisation \inserted{(i)}, car $X\to S'$ sera \unsure{qu. cpt|quotient}, donc couvrant pour fpqc donc $S'=X/R$, où $R=X\times_S X$.
Reste à prouver (ii) $\Rightarrow$ \inserted{le faisceau} $X/R$ est représentable par un $S'$ de \inserted{\unsure{loc.|prés.}} type fini sur $S$. \deleted{(N.B. alors, cela quitte à \unsure{part} sur $S$ car $X\to S'$ est}

[N.B. Alors $X\to S'$ est couvrant pour \unsure{fppc|fpqc|fppf}, comme après ch. \unsure{mt} de base $X\to S'$ il donne $X\times_S X=R\longrightarrow X$ qui est fid plat \unsure{par hypoth|qu. cpt}, il est fid plat qu. cpt, lui-$=$, ce qui implique \inserted{aussi} que $S'$ est \inserted{qu. cpt donc} de t.f. \deleted{\unsure{donc qu.|de présentation}} sur $S$].

Or \deleted{\unsure{comme|le foncteur} $X/R\to S$} le foncteur représenté par $X/R\hookrightarrow S$ est \unsure{ici|i.e.} le suivant :
\[
F(\deleted{\unsure{S}}T)=\begin{cases}
\emptyset&\text{si } X_T\ \text{non fid plat sur }T\ \inserted{[\Longleftrightarrow X_T\to T\ \text{couvre fppc}]}\\
\{\emptyset\}&\text{sinon}
\end{cases}
\]
D'ailleurs on est ramené au cas $S$ affine, \uuline{noethérien}. \deleted{et} Alors on applique le Th, qui nous dit qu'il suffit de vérifier \deleted{que} (L). Or sous la condition de cette condition, $X_T\to T$ est couvrant [car il existe un pt $x$ de la fibre spéciale, et $\unsure{\mathrm{Sp}|\mathrm{Spec}}\,\mathcal O_{X,x}$ \unsure{--|$\to$} $\mathrm{Sp}\,\mathcal O_{T,t_0}$ est fid plat] donc on gagne !
% archive sheet number: 12
